% !TEX program = xelatex
% =============================================================================
% resume · 中文一页专业简历（经典专业版式；示例：埃隆·马斯克，公开履历事实）
%
% 版式（中文招聘市场最常见的专业板式，全黑无灰）：
%       · 头部：左侧居中姓名区（姓名 / 英文名 / 定位语 / 联系方式），
%         右上角 2.5cm×3.5cm 证件照框位；
%       · 节标题黑体放大 + 细黑通栏线（头部与正文之间不再加线）；
%       · 条目一行三列：公司居左、职位居中、时间居右（年.月）；
%       · 要点用中文编号「1、2、」，关键数据与关键词用黑体强调。
% 字体（宋黑搭配，中文印刷排版的标准做法）：
%       · 正文宋体（ctex 按操作系统自动选择），标题与强调用黑体——
%         中文的"加粗"本质是字体切换（宋→黑），比伪粗更专业；
%       · 西文与数字 Source Sans Pro（TeX Live 自带），与黑体观感呼应。
% LaTeX 功力体现在细节：
%       · 时间列 OpenType 等宽数字（Numbers=Monospaced），右缘逐位对齐；
%       · microtype 光学边缘；xeCJK 中西文自动间距与全角标点挤压；
%       · 英文名 LetterSpace 字距；全文禁断词 + \emergencystretch 零溢出。
% 示例：以公众人物公开履历为样例内容（类比 neurips-paper 复刻公开论文），
%       联系方式为占位符，页脚有样例声明。
% 编译：bash scripts/compile.sh templates/resume/main.tex --preview
% =============================================================================
\documentclass[zihao=5, linespread=1.3]{ctexart}

% --- 版面 --------------------------------------------------------------------
\usepackage[a4paper, top=1.4cm, bottom=1.3cm, left=1.7cm, right=1.7cm]{geometry}
\pagestyle{empty}
\setlength{\parindent}{0pt}

% --- 字体与微排版 ------------------------------------------------------------
\usepackage[default]{sourcesanspro}   % 西文与数字：无衬线，与黑体呼应
\usepackage{microtype}                % 光学边缘（XeTeX 下启用 protrusion）
\usepackage{graphicx}                 % 证件照
\usepackage{enumitem}
\setlist[enumerate]{nosep, leftmargin=1.7em, itemsep=0.35ex, label=\arabic*、}
\setlist[itemize]{nosep, leftmargin=1.2em, itemsep=0.35ex,
                  label=\raisebox{0.18ex}{\small\textbullet}}
\hyphenpenalty=10000 \exhyphenpenalty=10000
\emergencystretch=8pt
\usepackage[hidelinks]{hyperref}
\hypersetup{pdftitle={埃隆·马斯克 — 简历（模板示例）}}

% --- 版式宏 ------------------------------------------------------------------
% 强调：中文排版惯例，正文宋体内用黑体（而非伪粗）标出关键词与数据
\newcommand{\hl}[1]{{\heiti\bfseries #1}}
% 时间样式：加粗、等宽数字（右缘逐位对齐）
\newcommand{\datefmt}[1]{{\bfseries\addfontfeature{Numbers=Monospaced}#1}}
% 节标题：黑体放大 + 细黑通栏线
\newcommand{\rsection}[1]{%
  \vspace{1.8ex}%
  {\Large\heiti #1}\par
  \nointerlineskip\vspace{0.6ex}%
  \hrule height 0.6pt\par
  \nointerlineskip\vspace{1.1ex}}
% 主条目（一行三列）：\entry{公司/学校}{职位/专业}{时间}
% 公司居左、职位居中、时间居右；三段零宽锚点互不挤压，居中列绝对居中
\newcommand{\entry}[3]{%
  \par\noindent
  \makebox[0pt][l]{\heiti\bfseries\large #1}\hfill
  \makebox[0pt][c]{\heiti #2}\hfill
  \makebox[0pt][r]{\datefmt{#3}}\par
  \vspace{0.45ex}}
% 单行条目：\entryline{公司}{一句话职责或结果}{时间}
\newcommand{\entryline}[3]{%
  \par\noindent
  {\heiti\bfseries #1}\ \ {#2}\hfill \datefmt{#3}\par\vspace{0.5ex}}

\begin{document}

% --- 头部：左侧姓名区，右上角证件照框位 --------------------------------------
\noindent
\begin{minipage}[c]{\dimexpr\textwidth-3.0cm\relax}
  {\centering
    {\zihao{2}\heiti\bfseries 埃隆·马斯克}\\[1.0ex]
    {\small\addfontfeature{LetterSpace=16}ELON MUSK}\\[1.1ex]
    企业家 · 工程师\\[0.55ex]
    美国得克萨斯州奥斯汀\quad|\quad
    \href{mailto:elon@example.com}{elon@example.com}\quad|\quad
    \href{https://x.com/elonmusk}{x.com/elonmusk}\par}
\end{minipage}\hfill
\begin{minipage}[c]{2.7cm}
  \raggedleft
  \setlength{\fboxrule}{0.5pt}\setlength{\fboxsep}{0pt}%
  % 证件照框位：换成真实照片时用下一行注释替换 \fbox{...} 整块
  % \includegraphics[width=2.5cm, height=3.5cm]{photo.jpg}
  \fbox{\parbox[c][3.5cm][c]{2.5cm}{\centering 证件照\\[0.3ex]
        {\scriptsize 2.5cm × 3.5cm}}}
\end{minipage}
\par\vspace{0.6ex}

% --- 个人简介 ----------------------------------------------------------------
\rsection{个人简介}
连续创业者，在支付、电动汽车、商业航天与人工智能领域先后创办多家定义行业的
公司：带领 SpaceX 从十余人的初创团队成长为\hl{全球发射市场的主导者}，带领
特斯拉从濒临破产走到\hl{全球市值最高的汽车公司}，长期以极限工程节奏交付
「首创级」产品。

% --- 工作经历 ----------------------------------------------------------------
\rsection{工作经历}

\entry{SpaceX}{创始人 · CEO 兼首席工程师}{2002.05 -- 至今}
\begin{enumerate}
  \item 研制 Falcon 1，成为\hl{全球首枚}由私营公司研制并成功入轨的液体燃料
        火箭（2008）；随后将 Falcon 9 打造为\hl{全球发射次数最多}的运载火箭；
  \item 开创火箭复用时代：\hl{2015 年}实现轨道级助推器首次动力回收着陆，
        \hl{2017 年}实现首次复飞，发射成本降低\hl{一个数量级}；
  \item \hl{首家}将宇航员送入国际空间站的私营公司（载人龙飞船，2020）；
        运营\hl{全球最大}卫星星座 Starlink 与完全可复用的 Starship 计划。
\end{enumerate}
\vspace{0.5ex}

\entry{特斯拉 Tesla, Inc.}{CEO 兼产品架构师}{2004.02 -- 至今}
\begin{enumerate}
  \item 2004 年领投 A 轮并出任董事长，\hl{2008 年}金融危机中接任 CEO，
        带领公司走出破产边缘；
  \item 相继推出 Roadster、Model S/X/3/Y 与 Cybertruck，成长为\hl{全球市值
        最高}的汽车制造商；
  \item 建成全球超级工厂（Gigafactory）制造网络与储能业务
        （Powerwall/Megapack），主导 \hl{Autopilot/FSD} 自动驾驶项目。
\end{enumerate}
\vspace{0.5ex}

\entry{xAI}{创始人}{2023.03 -- 至今}
\begin{enumerate}
  \item 前沿 AI 实验室，研发 \hl{Grok} 大模型系列；完成\hl{全球规模领先}的
        私募融资。
\end{enumerate}
\vspace{0.5ex}

\entryline{X 公司}{以 \hl{440 亿美元}收购 Twitter 并改组为 X。}{2022.10 -- 至今}
\entryline{Neuralink}{联合创始人——脑机接口，\hl{2024 年首例}人体植入。}{2016.07 -- 至今}
\entryline{The Boring Company}{创始人——建成拉斯维加斯会展中心环线。}{2016.12 -- 至今}
\entryline{OpenAI}{联合创始人兼联席主席（非营利研究阶段）。}{2015.12 -- 2018.02}
\entryline{X.com / PayPal}{创始人——合并后成为 PayPal，\hl{15 亿美元}售予 eBay。}{1999.03 -- 2002.10}
\entryline{Zip2}{联合创始人——城市指南软件，\hl{3.07 亿美元}售予康柏。}{1995.11 -- 1999.02}

% --- 教育背景 ----------------------------------------------------------------
\rsection{教育背景}
\entry{宾夕法尼亚大学}{物理学学士 · 经济学学士（沃顿商学院）}{1992 -- 1997}
\vspace{0.5ex}
\entry{斯坦福大学}{应用物理学博士项目}{1995.09}
入学两天后退学，创办 Zip2。

% --- 精选成就 ----------------------------------------------------------------
\rsection{精选成就}
\begin{itemize}
  \item \hl{首枚}私营入轨火箭 · \hl{首家}载人对接国际空间站的私营公司 ·
        \hl{首次}轨道级助推器回收与复飞；
  \item 2021 年\hl{《时代》年度人物}；六家公司横跨支付、能源、交通、隧道、
        神经科技与 AI。
\end{itemize}

\vspace{0.8ex}
{\centering\scriptsize
模板示例 —— 履历内容取自公开报道，联系方式为占位符。\par}

\end{document}
